\documentclass{article}
\usepackage{amsmath,amssymb}
\usepackage{xurl}
\usepackage[hidelinks]{hyperref}
% Shared commands for notes in docs/paper-gaps/.

\DeclareUnicodeCharacter{2080}{\ifmmode{}_{0}\else\textsubscript{0}\fi}
\DeclareUnicodeCharacter{2081}{\ifmmode{}_{1}\else\textsubscript{1}\fi}
\DeclareUnicodeCharacter{2082}{\ifmmode{}_{2}\else\textsubscript{2}\fi}
\DeclareUnicodeCharacter{2099}{\ifmmode{}_{n}\else\textsubscript{n}\fi}

\newcommand{\Lean}{\textsc{Lean}}
\ExplSyntaxOn
\NewDocumentCommand{\leanid}{m}{
  \ifmmode
    \textsf{\detokenize{#1}}
  \else
    \group_begin:
    \urlstyle{sf}
    \tl_set:Nn \l_tmpa_tl {#1}
    \tl_replace_all:Nnn \l_tmpa_tl {\_} {_}
    \exp_args:NV \path \l_tmpa_tl
    \group_end:
  \fi
}
\ExplSyntaxOff
\providecommand{\tr}{\operatorname{tr}}
\newcommand{\tnleanIssueBase}{https://github.com/LionSR/TNLean/issues/}
\newcommand{\tnleanPrBase}{https://github.com/LionSR/TNLean/pull/}
\newcommand{\tnleanDocBase}{https://sirui-lu.com/TNLean/docs/}
\newcommand{\ghissue}[1]{\href{\tnleanIssueBase#1}{GitHub issue \##1}}
\newcommand{\ghpr}[1]{\href{\tnleanPrBase#1}{PR \##1}}
\newcommand{\doclink}[1]{\href{\tnleanDocBase#1.html}{\path{#1}}}

% Dirac notation and the identity operator, as in blueprint/src/macros/common.tex.
% \providecommand keeps note-local definitions authoritative.
\usepackage{dsfont}
\providecommand{\ket}[1]{|#1\rangle}
\providecommand{\bra}[1]{\langle #1|}
\providecommand{\braket}[2]{\langle #1 | #2 \rangle}
\providecommand{\ketbra}[2]{|#1\rangle\!\langle #2|}
\providecommand{\Id}{\mathds{1}}
\providecommand{\one}{\mathds{1}}
\providecommand{\C}{\mathbb{C}}
\providecommand{\MN}[1]{M_{#1}(\C)}
\providecommand{\MD}{M_D(\C)}

% External referencing for paper sources.  The build helper writes sanitized
% auxiliary files under build/paper-aux/ and excludes duplicated source labels.
\usepackage{xr}

% Import a generated paper auxiliary file with a caller-supplied label prefix.
% The candidate paths support builds from docs/paper-gaps/, the repository root,
% and build/paper-aux/ itself.
\newcommand{\importpaperaux}[2]{%
  \IfFileExists{../../build/paper-aux/#2.aux}{%
    \externaldocument[#1]{../../build/paper-aux/#2}%
  }{%
    \IfFileExists{build/paper-aux/#2.aux}{%
      \externaldocument[#1]{build/paper-aux/#2}%
    }{%
      \IfFileExists{#2.aux}{%
        \externaldocument[#1]{#2}%
      }{}%
    }%
  }%
}

% General external reference macros
\newcommand{\paperref}[2]{\ref{#1-#2}}
\newcommand{\papereqref}[2]{(\ref{#1-#2})}

% CPSV16 source (1606.00608 / MPDO-22-12-17-2.tex) external document and shortcuts
\importpaperaux{cpsv16-}{1606.00608/MPDO-22-12-17-2-tnlean}
\newcommand{\cpsvref}[1]{\paperref{cpsv16}{#1}}
\newcommand{\cpsveqref}[1]{\papereqref{cpsv16}{#1}}

% Verdict marker: \gapnote{<kind>}{<status>}. Kind is the policy
% classification (clarification, local-correction, scope-restriction,
% unfaithful, false-source, open-gap); status is open, wip (elimination
% underway), resolved, or historical. Machine-read by the site index;
% prints a verdict line.
\newcommand{\gapnote}[2]{\par\noindent\textbf{Verdict:} #1 (#2).\par}

\title{GHZ Form of the Fixed-Point State: Local Orthogonality as a Hypothesis}
\author{}
\date{2026-09-26}
\begin{document}
\maketitle
\gapnote{scope-restriction}{open}

\section{What the source states}
Let $A^i=\bigoplus_{j=1}^b\operatorname{diag}(\mu_{j,1},\dots,\mu_{j,m_j})
    \otimes A_j^i$ be a basis of normal tensors, let $\ket{\omega_j}$ be the
fixed-point pair of the normal block $A_j$, embedded in the bond space
$\C^D\otimes\C^D$ of $A$, and put
\begin{align}
    \ket{\Omega'}
     & =\sum_{j=1}^b\alpha_j\bigotimes_{k=1}^{M}\ket{\omega_j}_{R_kL_{k+1}}.
    \label{eq:mswc24_fixed_point_state}
\end{align}
After eq.~(S7) of~\cite{Malz2024Preparation} the pairs are asserted to be
locally orthogonal,
\begin{align}
    \braket{\omega_j}{\omega_{j'}} & =\delta_{jj'},
    \label{eq:mswc24_local_orthogonality}
\end{align}
with a reference to the canonical-form theory of matrix product density
operators. The paragraph after eq.~(19) then concludes that
$W\colon\ket{j}\mapsto\ket{\omega_j}$ is an isometry and that
$\ket{\Omega'}=W^{\otimes M}\ket{\chi_M}$ with
$\ket{\chi_M}=\sum_j\alpha_j\ket{j}^{\otimes M}$. In the source,
(\ref{eq:mswc24_local_orthogonality}) is a property of these particular pairs,
not a hypothesis.

\section{The restriction adopted}
The formal statements are made for an arbitrary family
$\omega_1,\dots,\omega_b\in\C^D\otimes\C^D$ and arbitrary coefficients
$\alpha$.\footnote{Formal declarations:
    \leanid{MPSTensor.isIsometry_pairIsometry},
    \leanid{MPSTensor.nonNormalFixedPointState_eq_tensorPower_mulVec_ghzState}
    and \leanid{MPSTensor.nonNormalFixedPointState_norm_sq} in
    \doclink{TNLean/MPS/Preparation/NonNormalFixedPoint}.}
The identity $\ket{\Omega'}=W^{\otimes M}\ket{\chi_M}$ holds for every such
family, since
$W^{\otimes M}\ket{j}^{\otimes M}=\bigotimes_k\ket{\omega_j}$. The isometry of
$W$ and the normalization $\braket{\Omega'}{\Omega'}=1$ take
(\ref{eq:mswc24_local_orthogonality}) as a hypothesis. For the pairs of the
source they are therefore proved only once
(\ref{eq:mswc24_local_orthogonality}) is supplied.

One half of the hypothesis is available: if the components $\omega_{(a,c)}$ vanish
for $a\notin S$ and the components $\omega'_{(a,c)}$ vanish for $a\notin S'$,
where $S$ and $S'$ are disjoint sets of virtual indices, then $\braket{\omega}{\omega'}=0$.\footnote{Formal declaration:
    \leanid{MPSTensor.inner_pair_eq_zero_of_disjoint} in the same module.}
The embedded pairs of distinct blocks are supported in this way.

\section{Elimination}
Define the embedded pair of the $j$-th block in the bond space of $A$, show
that its right leg is supported on the virtual indices of that block, so that
distinct blocks give orthogonal pairs by the disjoint-support lemma, and show
$\braket{\omega_j}{\omega_j}=1$ from the normalization of the fixed-point pair
of a normal tensor in the gauge used by the source. The hypothesis
(\ref{eq:mswc24_local_orthogonality}) then holds for the source's pairs, and
the three statements above specialize to the source's claim without it.

\bibliographystyle{alpha}
\bibliography{references}
\end{document}
